\section{Mohammed and Charlemagne}

\begin{quotex}
Julius Evola's review of the book ``Mohammed and Charlemagne'' by Henri Pirenne, originally published in 1939.
\end{quotex}


People have often spoken of this book, especially since the recent release of the French edition: they attributed to him the merit of having upset the current idea of medieval civilization and its origins, thanks to his new and original views. This appreciation seems to us somewhat exaggerated. Certainly, we believe strongly that a history of the Middle Ages and its origins still needs to be written, given the numerous deformations and misunderstandings due to the rationalist and secularizing mentality of last century, which continue to prevail today on this subject. But to tell the truth, it is not in that sense that the Belgian researcher brings us new lights, just the opposite: after having begun his studies in the atmosphere of the economist school, the source of the doctrine of historical materialism, Pirenne depicts a further enlargement of its horizons which still feels the effects of the fundamental limitations of that school. We therefore find again in him an inclination to consider as the determining historical cause what, in our opinion, can have value at most as an accidental cause, so that to emphasize the economic and incidental period, conferring on it the greatest weight in the genesis of the forms of civilization.

This is the fundamental thesis defended by Pirenne in this, his final, work. That admirable construction which was the Roman Empire had an essentially Mediterranean character. The Mediterranean was its base and center. The sea, in all the force of the name \emph{Mare nostrum}, carried and spread ideas, religions, merchandise; it is the spiritual and material crucible of the North and the South, of the West and the East, under the sign of Rome, while the outer provinces are only a buffer against the Barbarians. The totality of life is concentrated on the shores of the great Roman lake. That said, according to Pirenne, the Germanic invasions did not put an end either to the Mediterranean unity of the ancient world, nor to what one can consider to be the essence of Roman civilization, as it still was in the fifth century, that is, an epoch where the Emperor of the West no longer existed. Byzantium/Constantinople assumes the patrimony of Rome as the Mediterranean capital. In spite of the havoc and diverse losses, that transfer did not determine the appearance of new guiding principles, neither in the economic order nor in the social order, nor in the state of the language, nor in those of the institutions. What remained of civilization was Mediterranean; civilization was conserved along the shores of the sea and it is from there that the new phenomena arise. The establishment of the Germanic barbarians in the Western provinces of the Empire produced nothing original; they did not intend to destroy the Empire nor to exploit it. Far from despising it, they admired it. Starting from the time of their establishment, all the heroic and original traits of the barbarian character have disappeared, according to Pirenne, absorbed by the Roman and Byzantine lands. And again following Pirenne, it is a fable to say that the Germans brought their morality of the young people, that is, of the people in whose heart personal connections of fidelity won over those of respect towards the State and was opposed to the cosmopolitan corruption of ancient civilization. That is the opposite which was produced and was true: the barbarians masterly knew how to adapt themselves to the Byzantine corruption, even contributing to aggravating it. Their structures were secular, they did not furnish any spiritual support. Constantinople remains the equivalent of the fallen Roman center. In the year 600, the world did not take a qualitatively different appearance to what it had in the year 400, and that principally because of the survival of the Mediterranean character of that civilization. The rupture with ancient Tradition and the Middle Ages (as a new civilization) begins therefore later than what is commonly supposed.

It is not the barbarian invasions which constituted an instrument of action, but rather a new phenomenon: the appearance of Islam with its rapid and unexpected progression. The consequence was the separation of the East and the West and the end of Mediterranean unity. Some countries like Africa and Spain, which had continued to be part of western civilization gravitated from that moment around the orbit of Baghdad. Islam represents a new religion and a new, unassimilable, aggressive, and irreducible civilization. By becoming a Muslim lake, the western Mediterranean ceased to be the route of cultural and economic exchanges of the Roman-Byzantine world, and the West found itself blocked and constrained to live folded in on itself. It is thus that the axis of the life of the world was displaced from the Mediterranean towards the North: not because of the effectiveness of spiritual and racial factors, but essentially as the result of that contingency. That provoked the crisis of the Merovingian regime, up until the appearance of a new dynasty originating from the Germanic countries of the North, the Carolingian dynasty. The pope was allied to it, detaching himself from the Emperor of the East, who, absorbed in the battle against the Muslims, was no longer in a position to sustain and protect him. A Rome detached from Byzantine could not avoid looking north. It is thus that the Church prepared the new order of things. It had assumed the symbol of Rome, and its authority became greater than the State, not succeeding in maintaining a strong centralized administration, it let itself be absorbed by feudalism which came to take on a particular importance, for the same reason of the paralysis of Mediterranean commerce and exchanges.

Dominated by the Church and feudalism, Europe is going to take on a new appearance. The Middle Ages begins in a birth which, according to Pirenne, occupies the century between 650 and 750. It was a period of anarchy where the ancient tradition was lost and where new elements got the upper hand. The evolution ceased in the year 800 which sealed the definitive rupture of the West and the East, the former giving itself a new Roman Germanic Empire, now detached from the old Roman Empire prolonged by Byzantium. It constituted a new civilization and a new civil community whose symbol and instrument was Carolingian Europe, on the base of a Germanic element Romanized by the Church. Even when political unity disappeared, there subsisted an international unity of civilization which was followed up to the Renaissance, an essentially Carolingian imprint.

These are the basic theses of Pirenne's book. This quick and easy note suffices to let the reader take account of the part that the materialist mentality and an extrinsic consideration of history plays in these theses. And, as always in similar cases, what could have been a necessary cause comes to be considered as the sufficient cause and as the \emph{conditio sine qua non}, from which the aberrant deductions arise. ``It is rigorously true,'' says Pirenne, ``that without Mohammed, Charlemagne is inconceivable.'' We have here the image of physical determinism according to which, for example, a thrust contains exactly everything that can explain the movement to which it passively constrained the inert body which receives the brunt of it. Pirenne doesn't even seem to suspect the \emph{active} causes of history.

The thesis asserted that the Roman event has survived much longer than believed in the invasion and the influence of the barbarians, that it had not disappeared at the end of the Empire but that the Mediterranean unity had continued to live, under the sign of Byzantium, from the fifth to the eight century, before being finally swept by Islam, and certain flattering for us and something to please us.

But, as the translator correctly points out in the preface, in order to sustain that thesis, it is not necessary to look beyond the façade, that is, to stick to the social and administrative structure of the ancient world, because behind that beautiful relatively intact appearance is found the putrefaction or something else from which the seed of that putrefaction emanates. It is a fact that, in order to assure the survival of the Roman event, Pirenne deliberately confuses it with the Byzantine event. He believes totally seriously that the process of the Byzantification of the Mediterranean world, which is produced up to the time of the Arab invasion, was able to truly assume the validity of the longevity of the ancient heritage.

A heritage which, in fact, had already undergone a fundamental alteration with the Christianization of the Mediterranean world. In order to not take that alteration into account, the author is again constrained to limit his considerations to the exterior, hierarchical, and political part of the Church, neglecting the properly spiritual factor. In reality, the survival of the Mediterranean unity had only a material and economic character, circumstantial for the most part. And it is just for that reason that an equally contingent circumstance, like the Arab invasion and the control by the Muslims of the maritime routes of the Mediterranean, was able to easily break it. That is in fact elsewhere than where one finds the positive forces preparing the coming of the new age, in which something of the purity of the ancient world had to truly survive invisibly.

One must certainly recognize, with Pirenne, the necessity of moving the date of the beginning of the Middle Ages, considered as a civilization having its own appearance, from the instant of the fall of the Western Empire to that of the creation of the Holy Roman German Empire -- that is, to move it ahead at least three centuries. These three centuries don't have a well-defined appearance. One could qualify them as the centuries of the interregnum. They do not yet correspond to a specific type of civilization. It is no longer Antiquity, but not yet the new times. Pirenne is however incorrect in not presenting the forces in gestation which already exercise a muddled action, in order to take ascendancy afterwards. He can be agreeable to us to hear him affirm and to see him struggle to demonstrate that the Germanic races lost their characteristics during their long entry into the Roman world and were ``Romanized''. But from nothing he can take nothing. It is not necessary to forget that, for the author, Romanization has often, we repeat, the meaning of Byzantization, consisting of assuming the vices and corruption of Byzantine decadence. At the point that he considers that what happened under the dynasty of the new Romanized barbarian States in the western basin of the Mediterranean was the decadence of a decadence. The author emphasizes almost tendentiously that corrupted, or at least passive and negative, aspect in order to be able to note the purely formal survival of Antiquity, while he ignores, at least in appearance, what there was positive in it.

He purely political deduction of the birth of the Holy Roman Empire is totally insufficient: he can expound it only to some minds located under the narcosis of materialism applied to history. To underline a fundamental point, Pirenne insists, for example, on the purely secular and profane character of the States and dynasties characteristic of the Germanic races up to the end of the Merovingian period -- it is only with the Carolingians and by the intervention of the Church which itself had been pushed by considerations of propriety that the royal power had assumed a sacred character. A more inexact affirmation that it was founded on an ambivalence making him believe that secularism exists everywhere where there is no indirect consecration granted by a distinct organization, that is, by the Church. In the same perspective, it would be necessary to deny the sacred character to a great part of the theocratic powers of Antiquity, and today even to those of the Japanese sovereign. In all these cases, it is an effect of tradition that a dynasty given can be sacred in itself and for itself, in virtue of some enigmatic heredity or destiny. Now, Pirenne seems to not know that what had already appeared, at the time of the invasions, in numerous ``barbarian'' tribes: among the Goths, they called the kings \emph{Asen} (``divine heroes'') and \emph{Amals} (``celestials'') because they believed that in the difference from the other tribes (or \emph{Sippen}), their line drew its origin from the mystical \emph{Asgard} and that their blood was not of a human nature. Certainly, we find these traditions in a state of involution and degeneration in the course of the period that we have called the interregnum, but not at the point of preventing, under certain conditions, a resumption, an awakening, a galvanization. That is exactly what happened, thanks to the tradition already mentioned and to a certain heroic morality, when an historic conjuncture led to the formation of the Holy Roman Empire. Here, the Roman symbol was the common reference of a supernatural order proposed by Christianity, which acted, so to say, as a catalyst on the manifestation and the survival of spiritual and racial forces latent and diverted in the course of the preceding period of Byzantization. Without admitting that latent survival of forces which then had to blossom under the form of a civilization having many affinities with the first Aryan societies, including the primordial Roman society, one could then explain the conflict between the Empire and the Church, between the Roman Ghibelline symbol and the Roman symbol of the Guelph. Or then, in Pirenne's manner and of so many other myopic historians, one will be obliged to explain that conflict only as the function of temporal interests and ambitions, thus missing the true essence of the conflict, the essence that we have often had the chance to bring to light in several of our works.

What it comes down to, if up until Mohammed the positive element had been only the Roman-Byzantium component, how does he explain that once the Mediterranean unity was swept away, the besieged West did not elaborate a civilization of the same type? Pirenne says that the Carolingians had no concern with a situation created from all pieces, but they drew parts from fortuitous circumstances, and that the coup d'état which replaced the Merovingian dynasty, the only one which survived after the invasions, is explained only by the closing of the Mediterranean by the Saracens. Even if he could truly justify such a coup d'état, that would never explain the genesis of the new type of civilization by which the Middle Ages took its own appearance and which, in its great political-social structures, reproduced essentially the distinctive traits of all the Aryan civilizations of Antiquity.

On this last point, one must likewise note that Pirenne expresses, it would only be incidentally, the judgments which would not receive our agreement: in his opinion, the passage of an almost cosmopolitan economy of exchange, credit and commerce, favored as it was by the Mediterranean routes, to an essentially agrarian and feudal economy which had to leave its imprint on medieval society properly called, signifies regression and decadence. For us, it is on the contrary the civilization not pursuing mercantile objectives which represents a value. And in that case also, the emergency provoked by the irruption of Islam in the Mediterranean acted like a stimulant, but not as a determining cause: as it happens, an exterior contingency served only to compel something superior, which was latent, to become manifest and take over. Beginning with the Carolingian period, the traffic of merchandise and commerce, with the mobility of riches, became the prerogative of the Jews, and only much later, that of the new centers foreshadowing the democratic and oligarchic-bourgeois corruption, so that the healthiest and purest forces of Germanic-Roman Europe separated themselves and developed their possibilities in a normal direction characteristic of a heroic, differentiated, and spiritual civilization.


\flrightit{Posted on 2016-02-28 by Cologero}

